\section{Bộ 46}
\debai{Trong mặt phẳng tọa độ Oxy, cho điểm M(0;2) và hai đường thẳng d1:x+2y=0, d2:4x+3y=0.Viết phương trình đường tròn đi qua điểm M có tâm thuộc đường thẳng d1 và cắt d2 tại hai điểm A,B sao cho độ dài đoạn AB bằng $4\sqrt{3}$. Biết tâm đường tròn có tung độ dương.}



\debai{Giải hệ phương trình:
$\begin{cases}x^{3}+12y^{2}+x+2=8y^{3}+8y\\ \sqrt{x^{2}+8y^{3}}=5x-2y  \end{cases}$}


\debai{TÌm giá trị nhỏ nhất của biểu thức: 
$$S=\frac{3}{b+c-a}+\frac{4}{a+c-b}+\frac{5}{a+b-c}$$
Trong đó a,b,c là độ dài các cạnh của một tam giác thỏa mãn $2a+b=abc$}


\section{Bộ 47}
\debai{Trong mặt phẳng với hệ tọa độ Oxy cho tam giác ABC có trực tâm H(3;0) và trung điểm của BC là I(6;1). Đường thẳng AH có phương trình x+2y-3=0. GỌi D;E lần lượt là chân đường cao kẻ từ B,C của tam giác ABC. Xác định tọa độ đỉnh của tam giác ABC, biết đường thẳng DE có phương trình x-2=0 và điểm D có tùn độ dương.}

\debai{Giải hệ phương trình:
$\left\{\begin{matrix}
xy+2=y\sqrt{x^{2}+2}(1) &  & \\ 
y^{2}+2\left ( x+1 \right )\sqrt{x^{2}+2x+3}=2x^{2}-4x(2) &  & 
\end{matrix}\right.$}

\debai{CHo x;y;z là các số thực dương thỏa mãn: x+y-x=-1. Tìm GTLN của biểu thức:
$$P=\frac{x^{3}y^{3}}{\left ( x+yz \right )\left ( y+xz \right )\left ( x+xy \right )^{2}}$$}


\section{Bộ 48}
\debai{GIải bất phương trình: $2x+5> \sqrt{2-x}\left ( \sqrt{x-1}+\sqrt{3x+4} \right )$}


\debai{Trong mặt phẳng tọa độ Oxy cho hình thang ABCD có $\angle BAD=\angle ADC=90^{\circ}$, AB=AD=2; DC=4, đỉnh C nằm trên đường thẳng d:2x-y+2=0. Điểm M nằm trên cạnh AD sao cho AM=2MD và đường thẳng BM có phương trình là 3x-2y+2=0. TÌm tọa độ C.}

\debai{Cho a;b;c là các số thực dương thỏa mãn: $3\left ( a^{2}+b^{2}+c^{2} \right )=1$. TÌm GTNN của biểu thức:
$$Q=\sqrt{a^{2}+b^{2}+\frac{1}{b^{2}}+\frac{1}{c^{2}}}+\sqrt{b^{2}+c^{2}+\frac{1}{c^{2}}+\frac{1}{a^{2}}}+\sqrt{c^{2}+a^{2}+\frac{1}{a^{2}}+\frac{1}{b^{2}}}$$}


\section{Bộ 49}
\debai{Giải hệ phương trình:
$\left\{\begin{matrix}
\left ( \sqrt{y}+1 \right )^{2}+\frac{y^{2}}{x}=y^{2}+2\sqrt{x-2} \hspace{2mm}(1)&  & \\ 
x+\frac{x-1}{y}+\frac{y}{x}=y^{2}+y \hspace{2mm}(2) &  & 
\end{matrix}\right.$}


\debai{Trong mặt phẳng tọa độ Oxy, họi H(3;-2),I(8;11),K(4;-1) lần lượt là trực tâm, tâm đường tròn ngoại tiếp, chân đường cao vẽ từ A của tam giác ABC. Tìm tọa độ A;B;C.}


\debai{Cho hai số thực x;y thỏa mãn điều kiện $x^{4}+16y^{4}+2\left ( 2xy-5 \right )^{2}=41$. TÌm GTNN-GTLN của biểu thức: $$P=xy-\frac{3}{x^{2}+4xy^{2}+3}$$}


\section{Bộ 50}
\debai{Giải hệ phương trình
$\left\{\begin{matrix}
x^{2}y+x^{2}+1=2x\sqrt{x^{2}y+2}\hspace{5mm} (1)&  & \\ 
y^{3}\left ( x^{6}-1 \right )+3y\left ( x^{2}-2 \right )+3y^{2}+4=0 \hspace{5mm}(2)&  & 
\end{matrix}\right.$}

\debai{Trong mặt phẳng hệ tọa độ Oxy, cho điểm E(3;4) đường thẳng $d:2+y-1=0$ và đường tròn $(C):x^{2}+y^{2}+4x-2y-4=0$. Gọi M là điểm thuộc đường thẳng d và nằm ngoài đường tròn (C). Từ M kẻ tiếp tuyến MA, MB đến đường tròn (C) (A, B là các tiếp điểm). Gọi (E) là đường tròn tâm E và tiếp xúc cới đường thẳng AB. Tìm tọa độ điểm M sao cho đường tròn (E) có chu vi lớn nhất.}

\debai{Cho x;y;z là các số thực thỏa mãn: $-1-2\sqrt{2}< x< -1+2\sqrt{2};y> 0;z> 0;x+y+z=-1$. Tìm GTNN của biểu thức:
$$P=\frac{1}{\left ( x+y \right )^{2}}+\frac{1}{\left ( x+z \right )^{2}}+\frac{1}{8-\left ( y+z \right )^{2}}$$}

\section{Bộ 51}
\debai{Giải hệ phương trình:
$\left\{\begin{matrix}
\sqrt{x-1}+\sqrt{y-1}=2 &  & \\ 
\sqrt{x+2}+\sqrt{y+2}=4 &  & 
\end{matrix}\right.$}

\debai{Trong mặt phẳng tọa độ $Oxy$ cho tam giác $ABC$ có $AB=3AC$. ĐƯờng phân giác trong của góc $\widehat{BAC}$ có phương trình $x-y=0$. Đường cao $BH$ có phương trình: 3x+y-16=0. Hãy xác định tọa độ $A,B,C$ biết $AB$ đi qua điểm $M(4;10)$}

\debai{Xét 3 số thực x;y;z thỏa mãn $x^{3}+y^{3}+z^{3}-3xyz=1$. TÌm GTNN của biểu thức: $$P=x^{2}+y^{2}+z^{2}$$}

\section{Bộ 52}
\debai{Giải phương trình:
$2x^{3}+9x^{2}-6x(1+2\sqrt{6x-1})+2\sqrt{6x-1}+8=0$}

\debai{Trong mặt phẳng hệ tọa độ Oxy cho tam giác ABC vuông tại A, có B(-2;1) và C(8;1). Đường tròn nội tiếp tam giác ABC có bán kính $r=3\sqrt{5}-5$. Tìm tọa độ tâm đường tròn nội tiếp tam giác ABC biết $y_{I}>0$.}


\debai{Cho 2 số thực dương tùy ý a;b;c. Tìm GTNN của biểu thức:

\begin{center}
$P=\frac{\sqrt{a^{3}c}}{2\sqrt{b^{3}a}+3bc}+\frac{\sqrt{b^{3}a}}{2\sqrt{c^{3}b}+3ca}+\frac{\sqrt{c^{3}b}}{2\sqrt{a^{3}c}+3ab}$
\end{center}}

\section{Bộ 53}
\debai{Giải hệ phương trình trên tập số thực:
$\left\{\begin{matrix}
2\sqrt{x^{2}+5}=2\sqrt{2y}+x^{2}\hspace{7mm}(1) &  & \\ 
x+3\sqrt{xy+x-y^{2}-y}=5y+4\hspace{7mm}(2) &  & 
\end{matrix}\right.$
}

\debai{Trong mặt phẳng tọa độ Oxy cho tứ giác ABCD nội tiếp đường tròn AC. Điểm M(3;-1) là trung điểm của BD, C(4;-2). Điểm N(-1;-3) nằm trên đường thẳng đi qua B và vuông góc với AD. Đường thẳng AD đi qua P(1;3). Tìm tọa độ A;B;D}

\debai{Cho x là số thực thuộc đoạn $\left [ -1;\frac{5}{4} \right ]$. TÌm GTLN, GTNN của biểu thức:
\begin{center}
$P=\frac{\sqrt{5-4x}-\sqrt{1+x}}{\sqrt{5-4x}+2\sqrt{1+x}+6}$
\end{center}}

\section{Bộ 54}
\debai{Giải hệ phương trình:
$\left\{\begin{matrix}
x^{2}+6y-4=\sqrt{2(1-y)(x^{3}+1)}\hspace{5mm}(1) &  & \\ 
(3-x)\sqrt{2-x}-2y\sqrt{2y-1}=0\hspace{5mm}(2) &  & 
\end{matrix}\right.$
}

\debai{Trong mặt phẳng tọa độ $Oxy$ cho hình chữ nhật $ABCD$ có diện tích bằng 12, $\left(\frac{9}{2};\frac{3}{2}\right)$ là tâm hình chữ nhật và $M(3;0)$ là trung điểm của $AD$. Tìm tọa độ các đỉnh của hình chữ nhật biết $y_{D}<0$}

\debai{Cho các số thực dương a;b;c thỏa mãn: $a^{2}+b^{2}+c^{2}=1$. TÌm GTNN của biểu thức:$$P=\frac{1}{\sqrt{a^{2}+ab}}+\frac{1}{\sqrt{b^{2}+ab}}+\frac{2\sqrt{3}}{1+c}$$}

\section{Bộ 55}
\debai{Giải bất phương trình: $\sqrt{4x+1}+\sqrt{6x+4}\geq 2x^{2}-2x+3$}

\debai{Trong mặt phẳng Oxy cho tam giác ABC có A(1;5), đường phân giác trong góc A có phương trình x-1=0, tâm đường tròn ngoại tiếp tam giác ABC là I($\frac{-3}{2}$;0) và điểm M(10;2) thuộc đường thẳng BC. Tìm tọa độ B;C}

\debai{Cho a;b;c là 3 số dương thỏa mãn $s^{2}+b^{2}+c^{2}=1$. TÌm GTNN của biểu thức: $$P=\frac{1}{a^{4}+a^{2}b^{2}}+\frac{1}{b^{4}+a^{2}b^{2}}+\frac{32}{(1+c)^{3}}$$}
